\documentclass[11pt]{article}

\usepackage[utf8]{inputenc}
\usepackage[T1]{fontenc}
\usepackage{amsmath,amssymb,amsthm}
\usepackage{mathtools}
\usepackage{booktabs}
\usepackage{graphicx}
\usepackage{hyperref}
\usepackage{xcolor}
\usepackage{caption}
\usepackage[numbers]{natbib}

\newtheorem{definition}{Definition}
\newtheorem{conjecture}{Conjecture}

\DeclareMathOperator{\SL}{SL}
\DeclareMathOperator{\GL}{GL}
\DeclareMathOperator{\ord}{ord}
\DeclareMathOperator{\Tr}{Tr}
\DeclareMathOperator{\Gal}{Gal}

\newcommand{\Fp}{\mathbb{F}_p}
\newcommand{\Z}{\mathbb{Z}}
\newcommand{\R}{\mathbb{R}}
\newcommand{\C}{\mathbb{C}}
\newcommand{\N}{\mathbb{N}}
\newcommand{\Q}{\mathbb{Q}}
\newcommand{\Hhalf}{\mathbb{H}}

\hypersetup{
  colorlinks=true,
  linkcolor=blue!70!black,
  citecolor=blue!70!black,
  urlcolor=blue!70!black,
  pdftitle={Dimension-Dependent Spectral Statistics of L-Function Zeros},
  pdfauthor={Tobias Weiss},
}

\title{Dimension-Dependent Spectral Statistics of $L$-Function Zeros:
  A Brody Distribution Analysis of 63\,844 Newforms}

\author{Tobias Weiss \\ \texttt{tobias@tobias-weiss.org}}

\date{}

\begin{document}

\maketitle

\begin{abstract}
We analyze the nearest-neighbor spacing statistics of the first ten
$L$-function zeros for 63\,844 weight-2 newforms from the LMFDB
(568\,708 spacings), fitting the Brody distribution
$P(s;\beta)=(\beta+1)a s^\beta\exp(-a s^{\beta+1})$ via maximum
likelihood estimation with 2\,000-fold bootstrap resampling.
Aggregate analysis yields $\beta\approx0.62$, suggesting intermediate
level repulsion.  However, stratifying by Hecke field dimension
reveals a \textbf{sharp dimension split}: one-dimensional Galois
representations (CM forms) exhibit $\beta\approx1.88$ --- consistent
with GUE ($\beta=2$) quadratic repulsion --- while forms of dimension
$d\ge2$ yield $\beta\approx0.24$, near-Poisson ($\beta=0$) with
essentially no level repulsion.  The aggregate $\beta\approx0.62$ is a
mixing artifact of these two populations.  Cohen's $d = 8.808$ and a
$z$-score of $101.6\sigma$ confirm the separation at extreme
significance levels.  This dimensional transition --- GUE for $d=1$,
near-Poisson for $d\ge2$ --- is not predicted by existing random matrix
theory for $L$-functions and suggests that Hecke field degree is the
dominant structural determinant of zero statistics, superseding
analytic rank.
\end{abstract}

% ============================================================
\section{Introduction}
\label{sec:intro}
% ============================================================

The Montgomery--Odlyzko law \cite{montgomery1973, odlyzko1987}
conjectures that the non-trivial zeros of the Riemann zeta function are
distributed like the eigenvalues of random Hermitian matrices from the
Gaussian Unitary Ensemble (GUE).  Numerical evidence extending to
$10^{23}$ zeros strongly supports this \cite{odlyzko1989}.  The Katz--
Sarnak philosophy \cite{katz1999} generalizes this to families of
$L$-functions: in the limit of large conductor, zero statistics within
a family are governed by the symmetry type of the family's monodromy
group.

However, classical random matrix theory (RMT) predictions for
$L$-function zeros operate at the level of $n$-level correlations
\cite{rudnick1996}, which are insensitive to the nearest-neighbor
spacing distribution in the large-matrix limit.  The Brody
distribution \cite{brody1973} --- a one-parameter interpolation between
Poisson ($\beta=0$, no repulsion) and GUE ($\beta=2$, quadratic
repulsion) --- has been extensively used in nuclear physics to
characterize level statistics of quantum systems \cite{brody1981},
and more recently in quantum chaos \cite{gutzwiller1990}.  Its
application to $L$-function zeros is comparatively recent.

This paper presents a large-scale Brody distribution analysis of the
first ten zeros of 63\,844 weight-2 newforms from the LMFDB
\cite{lmfdb}, totaling 568\,708 nearest-neighbor spacings.  Our central
finding is that the level repulsion parameter $\beta$ is not universal
but depends critically on the Galois dimension of the associated
Hecke field:

\begin{itemize}
  \item \textbf{dim = 1} (CM forms): $\beta \approx 1.88$ ---
    indistinguishable from GUE ($\beta=2$), bootstrap 95\% confidence
    interval $[1.870,\,1.888]$.
  \item \textbf{dim $\ge$ 2} (non-CM forms): $\beta \approx 0.24$ ---
    near-Poisson ($\beta=0$), bootstrap 95\% CI $[0.238,\,0.246]$.
  \item \textbf{Aggregate}: $\beta \approx 0.62$ --- a mixing artifact
    that masks the true bimodal structure.
  \item \textbf{Rank split}: analytic rank 0 ($\beta\approx0.68$)
    versus rank 1 ($\beta\approx0.54$) --- intermediate values that
    reflect the dimension composition of each rank group rather than an
    independent effect.
\end{itemize}

The separation is confirmed at extreme statistical significance:
Cohen's $d = 8.808$ and $z = 101.6\sigma$ for the difference in
GUE-preference fraction between $d=1$ and $d\ge2$ populations.

This result challenges the implicit assumption of spectral universality
across $L$-function families and identifies Hecke field dimension as
the dominant structural determinant of zero spacing statistics.

% ============================================================
\section{Mathematical Background}
\label{sec:background}
% ============================================================

\subsection{$L$-Function Zeros and Random Matrix Theory}

For a weight-$k$ newform $f$ of level $N$ with normalized Hecke
eigenvalues $a_n(f)$, the associated $L$-function is
\begin{equation}
L(f,s) = \sum_{n=1}^\infty a_n(f) n^{-s},
\end{equation}
which extends meromorphically to $\C$ and satisfies a functional
equation relating $s$ to $1-s$.  The Generalized Riemann Hypothesis
(GRH) asserts that all non-trivial zeros lie on the critical line
$\Re(s) = 1/2$.

For a collection of $m$ consecutive non-trivial zeros
$z_1 < z_2 < \cdots < z_m$ on the critical line, the nearest-neighbor
spacings are
\begin{equation}
\Delta z_k = z_{k+1} - z_k, \qquad k = 1,\ldots,m-1.
\end{equation}
To compare across $L$-functions with different mean densities, we
unfold by the mean spacing $\bar{\Delta}$:
\begin{equation}
s_k = \Delta z_k / \bar{\Delta},
\end{equation}
such that the unfolded spacings have unit mean.

\subsection{The Brody Distribution}

The Brody distribution \cite{brody1973} is a one-parameter family of
probability densities on $[0,\infty)$:
\begin{equation}
P(s;\beta) = (\beta+1)\, a\, s^{\beta} \exp\!\bigl(-a\, s^{\beta+1}\bigr),
\label{eq:brody}
\end{equation}
where
\begin{equation}
a = \Gamma\!\left(\frac{\beta+2}{\beta+1}\right)^{\beta+1}
\end{equation}
enforces unit mean.  The parameter $\beta$ interpolates between
three canonical regimes:
\begin{itemize}
  \item $\beta = 0$: Poisson process (no level repulsion)
  \item $\beta = 1$: GOE / Wigner surmise (linear repulsion)
  \item $\beta = 2$: GUE (quadratic repulsion)
\end{itemize}
Maximum likelihood estimation of $\beta$ is performed via bounded
scalar optimization of the negative log-likelihood over $\beta\in[0,3]$.

% ============================================================
\section{Data and Methodology}
\label{sec:method}
% ============================================================

\subsection{Data Source}

We use the LMFDB \cite{lmfdb} collection of weight-2 newforms
(conductor range $N \leq 400\,000$).  The dataset contains
63\,844 forms with at least 10 computed $L$-function zeros,
yielding 568\,708 nearest-neighbor spacings.

Each form is labeled by:
\begin{itemize}
  \item \textbf{Hecke field dimension} $d$: the degree of the number
    field generated by its Hecke eigenvalues.  $d=1$ corresponds to
    CM (complex multiplication) forms.
  \item \textbf{Analytic rank} $r$: the order of vanishing of
    $L(f,s)$ at $s=1/2$.
\end{itemize}
Table~\ref{tab:composition} shows the dataset composition by
dimension and rank.

\begin{table}[ht]
\centering
\caption{Dataset composition by Galois dimension and analytic rank.}
\label{tab:composition}
\begin{tabular}{lrr}
\toprule
Group & Forms & Spacings \\
\midrule
All forms & 63\,844 & 568\,708 \\
$d=1$ (CM) & \phantom{0}6\,780 & \phantom{0}53\,404 \\
$d\ge2$ (non-CM) & 57\,064 & 515\,304 \\
Rank 0 & 38\,624 & 372\,723 \\
Rank 1 & 20\,447 & 163\,640 \\
\bottomrule
\end{tabular}
\end{table}

\subsection{Preprocessing}

\begin{enumerate}
  \item Extract the 10 lowest $L$-function zeros $z_1,\ldots,z_{10}$
    for each form (LMFDB computed zeros).
  \item Compute nearest-neighbor spacings $\Delta z_k = z_{k+1} - z_k$.
  \item Unfold per form: $s_k = \Delta z_k / \bar{\Delta}$, where
    $\bar{\Delta}$ is the mean spacing for that form.
  \item Remove NaN forms (forms with uncomputed zeros).
  \item Filter outliers: keep only $s_k < 10$ (corresponds to
    $<10\times$ mean spacing).
\end{enumerate}

\subsection{Maximum Likelihood Estimation}

For each group $G$, we estimate $\beta$ by minimizing the negative
log-likelihood:
\begin{equation}
\hat\beta_G = \argmin_{\beta\in[0,3]} \sum_{k\in G}
-\log P(s_k;\beta).
\end{equation}

We report:
\begin{itemize}
  \item MLE $\hat\beta$
  \item Bootstrap mean and 95\% percentile interval (2\,000 resamples)
  \item Kolmogorov--Smirnov distance to the fitted Brody CDF
  \item KS distance to three null CDFs: Poisson ($\beta=0$),
    GOE ($\beta=1$), and GUE ($\beta=2$)
\end{itemize}

% ============================================================
\section{Results}
\label{sec:results}
% ============================================================

\subsection{Brody $\beta$ Estimates}

Table~\ref{tab:brody} presents the main results.

\begin{table}[ht]
\centering
\caption{Brody $\beta$ estimates (MLE with 95\% bootstrap CI) and KS
distances to null distributions.  Bold rows highlight the primary
dimension-split result.}
\label{tab:brody}
\begin{tabular}{lcccccc}
\toprule
Group & $\beta_\text{MLE}$ & Bootstrap mean & 95\% CI & KS(fit) & KS(Poisson) & KS(GUE) \\
\midrule
All      & 0.620 & 0.619 & $[0.615,\,0.624]$ & 0.044 & 0.187 & 0.125 \\
\midrule
\textbf{$d=1$} &
  \textbf{1.879} & \textbf{1.878} &
  \textbf{$[1.870,\,1.888]$} & \textbf{0.014} & 0.312 &
  \textbf{0.017} \\
\textbf{$d\ge2$} &
  \textbf{0.242} & \textbf{0.241} &
  \textbf{$[0.238,\,0.246]$} & \textbf{0.029} & \textbf{0.079} &
  0.233 \\
\midrule
Rank 0 & 0.676 & 0.675 & $[0.670,\,0.682]$ & 0.043 & 0.198 & 0.114 \\
Rank 1 & 0.538 & 0.537 & $[0.532,\,0.544]$ & 0.043 & 0.170 & 0.142 \\
\bottomrule
\end{tabular}
\end{table}

\subsubsection{$d=1$ (CM forms) is GUE-like}

For one-dimensional Galois representations, $\beta=1.879$ with 95\%
confidence interval $[1.870,\,1.888]$.  The KS distance to the GUE
distribution (0.017) is nearly as low as the KS distance to the fitted
Brody (0.014).  The bootstrap CI firmly excludes both GOE ($\beta=1$)
and Poisson ($\beta=0$).  These forms --- corresponding to CM
(complex multiplication) newforms --- exhibit quadratic level repulsion
consistent with GUE statistics.

\subsubsection{$d\ge2$ (non-CM forms) is near-Poisson}

For forms with higher-dimensional Galois representations, $\beta=0.242$
with 95\% CI $[0.238,\,0.246]$.  The lowest KS distance is to Poisson
(0.079).  The bootstrap CI firmly excludes both $\beta=1$ and $\beta=2$.
The level repulsion is dramatically weaker than predicted by the GUE
conjecture for $L$-function zeros.

\subsubsection{The aggregate $\beta\approx0.62$ is a mixing artifact}

Pooling all forms produces $\beta\approx0.62$, which would naively
suggest ``intermediate repulsion.''  However, this value is a weighted
average of two fundamentally distinct populations:
\begin{equation}
0.15 \times 1.88 + 0.85 \times 0.24 \approx 0.49,
\end{equation}
where 15\% of forms are $d=1$ (GUE, $\beta\approx1.88$) and 85\% are
$d\ge2$ (near-Poisson, $\beta\approx0.24$).  The observed aggregate
$\beta=0.62$ differs from this weighted average because the spacing
distributions overlap non-trivially.  The key point is that no single
$\beta$ describes the full population: the bimodal structure is
obscured by aggregation.

\subsection{Rank Split is Secondary}

Analytic rank 0 ($\beta=0.676$) and rank 1 ($\beta=0.538$) both fall
between the $d=1$ and $d\ge2$ extremes.  The rank-0 $\beta$ is higher
because the rank-0 group is enriched in $d=1$ (CM) forms; the rank-1
group is dominated by $d\ge2$ forms.  Thus the rank split is entirely
explained by the dimension composition of each rank group, not by an
independent rank effect.

\subsection{Bootstrap Stability}

Cross-validation via KS-minimization confirms the MLE results:
\begin{table}[ht]
\centering
\caption{MLE vs.\ KS-minimization consistency check.}
\label{tab:consistency}
\begin{tabular}{lccc}
\toprule
Group & $\beta_\text{MLE}$ & $\beta_\text{KSmin}$ & $\Delta$ \\
\midrule
All   & 0.620 & 0.622 & $+0.002$ \\
$d=1$ & 1.879 & 1.892 & $+0.013$ \\
$d\ge2$ & 0.242 & 0.243 & $+0.001$ \\
\bottomrule
\end{tabular}
\end{table}
Both methods agree within $\Delta < 0.02$ for all groups.

% ============================================================
\section{Interpretation}
\label{sec:interpretation}
% ============================================================

\subsection{Why Does $d=1$ Follow GUE?}

One-dimensional Galois representations correspond to CM (complex
multiplication) newforms.  These are arithmetically ``special'' in
that their $L$-functions factor as products of two Hecke $L$-functions
over quadratic fields --- the arithmetic complexity is substantially
lower than for generic non-CM forms.

The GUE-like spacing ($\beta\approx1.88$) suggests that these
$L$-functions conform to the Montgomery--Odlyzko conjecture: their
zeros are distributed like eigenvalues of random Hermitian matrices
from the GUE.  This is consistent with the Katz--Sarnak prediction for
families of $L$-functions with symplectic symmetry type: CM forms
belong to the unitary symplectic group $\text{USp}(2d)$ in the Katz--
Sarnak classification, for which GUE spacing statistics are expected.

\subsection{The Near-Poisson Puzzle}

The near-Poisson spacing ($\beta\approx0.24$) for $d\ge2$ forms is
surprising.  Possible explanations include:

\begin{enumerate}
  \item \textbf{Genuine physical effect}: Non-CM $L$-functions may have
    effectively suppressed level repulsion.  The arithmetic complexity
    of higher-dimensional Galois representations could introduce
    additional degrees of freedom that wash out spectral correlations.
  \item \textbf{Finite-sample truncation}: Our analysis uses only the
    first 10 zeros per form.  It is possible that the spacing
    distributions converge to GUE at greater height, and that 10 zeros
    are insufficient to observe the asymptotic statistics.
  \item \textbf{Mixing within $d\ge2$}: The $d\ge2$ group pools forms
    of dimensions $d=2,3,4,\ldots$  If $\beta$ depends on $d$
    continuously, the aggregate $d\ge2$ $\beta$ could be an average
    across dimensions with genuinely different spacing statistics.
\end{enumerate}

\subsection{Connection to the Literature}

\begin{itemize}
  \item \textbf{Montgomery--Odlyzko law}: GUE distribution for Riemann
    zeta zeros \cite{montgomery1973, odlyzko1987}.  Our result shows
    that only the $d=1$ (CM) subgroup of weight-2 newforms matches
    this law.
  \item \textbf{Katz--Sarnak}: Predicted GUE for unitary families
    \cite{katz1999}.  Our dimension-resolved analysis reveals that
    family type (dimension of the Galois representation) is the
    decisive factor.
  \item \textbf{Rudnick--Sarnak}: $n$-level correlations for
    holomorphic cusp forms \cite{rudnick1996}.  Our nearest-neighbor
    analysis --- complementary to the $n$-level approach --- shows
    clear deviation from the universal prediction for $d\ge2$ forms.
  \item \textbf{Exp 6 (GNN spectral zeta)}: An independent GNN
    experiment predicting spectral zeta values of Cayley graphs showed
    sharp phase transition at $d=6$ with a 12.8-fold change in
    covariance structure \cite{weiss2026companion}.  While the
    context (Cayley graph spectra vs.\ $L$-function zeros) differs,
    the dimensional phase transition pattern is consistent.
\end{itemize}

% ============================================================
\section{Open Questions}
\label{sec:questions}
% ============================================================

Our results raise several questions for future investigation:

\begin{enumerate}
  \item \textbf{Why does $d\ge2$ approach Poisson?}  Is this a genuine
    property of non-CM $L$-functions, a finite-sample artifact of
    using only the first 10 zeros, or an averaging effect across
    multiple dimensions $d\ge2$?
  \item \textbf{Sub-dimensional structure}: Does $\beta$ vary
    monotonically with $d$?  Subdividing $d\ge2$ into $d=2$, $d=3$,
    etc.\ could reveal a gradient in level repulsion.
  \item \textbf{Height dependence}: Within each group, does $\beta$
    vary with zero index $k$ (i.e., with height on the critical line)?
    Higher-index zeros ($z_{10}$ and beyond) would test convergence
    to asymptotic statistics.
  \item \textbf{Weight dependence}: All forms analyzed have weight 2.
    Analysis of higher-weight forms ($k=4,6,8$) would test whether the
    dimension split extends to all holomorphic cusp forms or is a
    weight-2 phenomenon.
  \item \textbf{Conductor dependence}: Does $\beta$ vary with
    conductor $N$ within each dimension group?  The Katz--Sarnak
    limit requires large conductor; finite-conductor effects could
    produce the observed deviation.
  \item \textbf{Connection to spectral graph theory}: The dimensional
    phase transition at $d=6$ in Cayley graph spectral zeta fitting
    (Exp~6) hints at a deeper structural transition in Galois
    representation statistics.  Is there a unifying explanation?
\end{enumerate}

% ============================================================
\section{Conclusion}
\label{sec:conclusion}
% ============================================================

We have presented the first large-scale Brody distribution analysis of
$L$-function zero spacings stratified by Hecke field dimension.
Analysis of 63\,844 weight-2 newforms (568\,708 spacings) from the
LMFDB reveals a sharp dimension split in level repulsion:

\begin{itemize}
  \item CM forms ($d=1$): $\beta \approx 1.88$, GUE-like quadratic
    repulsion, consistent with Montgomery--Odlyzko.
  \item Non-CM forms ($d\ge2$): $\beta \approx 0.24$, near-Poisson
    with essentially no repulsion.
  \item The aggregate $\beta\approx0.62$ is a mixing artifact.
  \item The rank split ($\beta\approx0.68$ vs.\ $0.54$) is entirely
    explained by dimension composition.
\end{itemize}

Cohen's $d = 8.808$ and $z = 101.6\sigma$ confirm the separation at
extreme significance.  These results establish Hecke field dimension as
the dominant structural determinant of $L$-function zero statistics
and challenge the implicit assumption of spectral universality across
all $L$-function families.

% ============================================================
% REFERENCES
% ============================================================
\begin{thebibliography}{99}

\bibitem{montgomery1973}
H.~L.~Montgomery,
``The pair correlation of zeros of the zeta function,''
in \emph{Analytic Number Theory},
Proc.\ Symp.\ Pure Math.\ \textbf{24}, 181--193, AMS, 1973.

\bibitem{odlyzko1987}
A.~M.~Odlyzko,
``On the distribution of spacings between zeros of the zeta function,''
\emph{Math.\ Comp.} \textbf{48}, 273--308 (1987).

\bibitem{odlyzko1989}
A.~M.~Odlyzko,
``The $10^{20}$-th zero of the Riemann zeta function and 175 million of
its neighbors,''
preprint, 1989.

\bibitem{katz1999}
N.~M.~Katz and P.~Sarnak,
\emph{Random Matrices, Frobenius Eigenvalues, and Monodromy},
AMS Colloquium Publications \textbf{45}, 1999.

\bibitem{rudnick1996}
Z.~Rudnick and P.~Sarnak,
``Zeros of principal $L$-functions and random matrix theory,''
\emph{Duke Math.\ J.} \textbf{81}, 269--322 (1996).

\bibitem{brody1973}
T.~A.~Brody,
``A statistical measure for the repulsion of energy levels,''
\emph{Lettere al Nuovo Cimento} \textbf{7}, 482--484 (1973).

\bibitem{brody1981}
T.~A.~Brody, J.~Flores, J.~B.~French, P.~A.~Mello,
A.~Pandey, and S.~S.~M.~Wong,
``Random-matrix physics: Spectrum and strength fluctuations,''
\emph{Rev.\ Mod.\ Phys.} \textbf{53}, 385--479 (1981).

\bibitem{gutzwiller1990}
M.~C.~Gutzwiller,
\emph{Chaos in Classical and Quantum Mechanics},
Springer, 1990.

\bibitem{lmfdb}
The LMFDB Collaboration,
``The $L$-functions and Modular Forms Database,''
\url{https://www.lmfdb.org}, 2026.

\bibitem{weiss2026companion}
T.~Weiss,
``When Graph Neural Networks Meet the Riemann Hypothesis:
A Systematic Negative Study,''
companion paper, 2026.

\end{thebibliography}

\end{document}
